\documentclass[../../../include/open-logic-section]{subfiles}

\begin{document}

\olfileid{int}{pty}{ter}

\olsection{证明项}

在 \Log{N1i} 中，我们用变元（例如 $!A^x$）标记假设，作为解除标签。
在 \Log{N2i} 中，每个相继式 $\Gamma \Sequent !A$ 左边的上下文~$\Gamma$ 中包含一个“变元‑公式”对的列表。
因此，在 \Log{N2i} 中，推导中的每个相继式不仅携带了到目前为止该推导推导出了什么（即 $!A$），
还携带了在每一步哪些假设是开放的（即 $\Gamma$）。
那么，如果我们在每个相继式中再加入 $!A$ 是\emph{如何}推导出来的信息，会怎样呢？
为此，我们可以定义一个\emph{带项标签的}系统 \Log{tN2i}，
其中相继式具有 $\Gamma \Sequent N : !A$ 的形式。
$\Gamma$ 和 $!A$ 的含义同前：$\Gamma$ 是形如 $x: !B$ 的配对集合，
$!A$ 则是以该相继式结尾的子推导表明可从~$\Gamma$ 中的假设推出的公式。
$N$ 将是某个项，用于编码以该相继式结尾的推导。
% part: intuitionistic-logic
% chapter: propositions-as-types
% section: proof-terms

我们分配给每个相继式的项是由变元 $x$, $y$,
……以及函数符号构建的，这些函数符号编码了到目前为止所执行的推理。每个推理规则都需要一个独立的函数符号。
对应引入规则的函数符号称为“构造子”，对应消去规则的则称为“析构子”。
每个函数符号所取的自变元个数，与其对应规则所具有的前提个数相同。
例如，对应 \Intro{\land} 和 \Elim{\land}[1] 的函数符号可以是 $c^\land$（二元）和 $d^\land_1$（一元）。
在 \Log{tN2i} 中，这些规则就变为：
\[
\Axiom$\Gamma_1 \fCenter N: !A$
\Axiom$\Gamma_2 \fCenter M: !B$
\RightLabel{\Intro{\land}}
\BinaryInf$\Gamma_1, \Gamma_2 \fCenter c^\land(N, m) : !A \land !B $
\DisplayProof
\quad
\Axiom$\Gamma \fCenter N: !A \land !B$
\RightLabel{\Elim{\land}[1]}
\UnaryInf$\Gamma \fCenter d^\land_1(N) : !A$
\DisplayProof
\]
如果一个推理规则解除了一个假设，也就是说，该规则有一个前提标明了带标签的公式~$x:!A$，
而该公式不在结论的上下文之中，那么我们必须在证明项中也标明这一点。
与此类规则相关联的函数符号会\emph{约束}相应的变元。
例如，\Intro{\lif} 规则从前提 $x:!A, \Gamma \Sequent !B$ 得到 $\Gamma \Sequent !A \lif !B$。
构造子 $c^{\lif}$ 只取一个自变元，但同时还必须指明变元 $x$ 被约束了。
这就是证明项表达带标签~$x$ 的假设已被解除的方式。
例如，如果前提的证明项是~$N$，我们可以将结论的证明项写作 $c^{\lif}([x]N)$，并将该规则写作
\[
\Axiom$x:!A, \Gamma \fCenter N: !B$
\RightLabel{\Intro{\lif}}
\UnaryInf$\Gamma \fCenter c^{\lif}([x]N) : !A \lif !B$
\DisplayProof
\]

为了能完全重建一个推导，我们还需要一些额外信息。
当一条规则的前提公式不能完全确定结论公式是什么时，我们就必须把这一信息添加到项中。
\Intro{\lif} 就属于这种情况，因此我们将写作 $c^{\lif}([x:!A]N)$。
\Intro{\lor} 和~\FalseInt 也属于这种情况，所以我们使用携带该信息的函数符号，例如
\[
\Axiom$\Gamma \fCenter N:!A$
\RightLabel{\Intro{\lor}[1]}
\UnaryInf$\Gamma \fCenter c^{\lor{!B}}_1:!A \lor !B$
\DisplayProof
\quad
\Axiom$\Gamma \fCenter N:\lfalse$
\RightLabel{\FalseInt}
\UnaryInf$\Gamma \fCenter d^\lfalse_{!A}(N):!A$
\DisplayProof
\]

基于这些想法，可以构造一个相继式风格的自然演绎系统，其中每个相继式的形式都是 $\Gamma \Sequent N : !A$，这里 $N$~是一个\emph{证明项}。

结果表明，这类证明项与所谓的\emph{类型$\lambda$演算}中的项有着密切的联系。
为了体现这一联系，我们将采用$\lambda$演算文献中使用的符号，而非上面那些通用的构造子与析构子符号。
例如，我们不写 $c^{\lif}([x:!A]N)$ 而写 $\lambd[\typeof{x}{!A}][N]$；用 $(NM)$ 代替 $d^{\lif}(N,M)$；
用 $\pair{N}{M}$ 代替 $c^\land(N,M)$；用 $\proj{i}{N}$ 代替~$d^\land_i(N)$。

\begin{defn}
  证明项由以下规则归纳生成：
  \begin{enumerate}
  \item 每个变元 $x$ 都是一个证明项（公理）。
  \item 若 $x$ 是一个变元，$!A$ 是一个公式，且 $N$ 是一个证明项，则 $\lambd[\typeof{x}{!A}][N]$ 也是一个证明项（\Intro\lif）。
  \item 若 $N$ 和 $M$ 是证明项，则 $(NM)$ 也是一个证明项（\Elim\lif）。
  \item 若 $N$ 和 $M$ 是证明项，则 $\pair{N}{M}$ 也是一个证明项（\Intro\land）。
  \item 若 $N$ 是一个证明项，则 $\proj{i}{N}$ 也是一个证明项，其中 $i$ 为 $1$ 或 $2$（\Elim\land[i]）。
  \item 若 $N$ 是一个证明项，且 $!A$ 是一个公式，则 $\inj[!A]{i}{!M}$ 也是一个证明项，其中 $i$ 为 $1$ 或 $2$（\Intro\lor[i]）。
  \item 若 $M$、$N_1$、$N_2$ 是证明项，且 $x_1$、$x_2$ 是变元，则 $\dcase{M}{x_1}{N_1}{x_2}{N_2}$ 也是一个证明项（\Elim\lor）。
  \item 若 $N$ 是一个证明项且 $!A$ 是一个公式，则 $\abort{!A}{N}$ 也是一个证明项（\FalseInt）。
  \end{enumerate}
\end{defn}

并非每个证明项都是某个推导的证明项。例如，
项 $\pair{N}{M}$ 编码了一个以~\Intro{\land} 推理结尾的证明，
因此它的结论公式是一个合取式~$!A \land !B$。
这个合取式不能成为~\Elim{\lif} 规则的主前提，
所以 $(\pair{N}{M}P)$ 不可能是一个正确证明的证明项。
只有那些遵循 \Log{tN2ip} 规则的证明项才是\emph{正确}的证明项。
这些规则见 \olref{tab:tN2ip}。

\subfile{rules-tN2}

\end{document}